% GENERATED FILE. Edit canonical lessons, then run npm run generate:books.
% canonical-source-hash: fnv1a64:c74f33f851acbcca
% canonical-lessons: HI-C378-avasyak, HI-C378-anivarya, HI-C378-keval, HI-C378-sivay, HI-C378-manakarna

\chapter{What Must Happen First}
\label{ch:hi-what-must-happen-first}
\hlchaptermodality{378}

\begin{chapteropening}
\textbf{By the end of this chapter:} \emph{I can say what is required before something can happen.}

Complete the last lesson of chapter 378: I can say what is required before something can happen.
\end{chapteropening}

\section[\texorpdfstring{\hi{आवश्यक} (āvaśyak)}{आवश्यक (āvaśyak)}]{\hi{आवश्यक} (āvaśyak) — essential, required (formal)}

\label{lesson:HI-C378-avasyak}



\begin{quote}
Before the new one: say the Hindi for a hardship, then the Hindi for to trap.
\end{quote}

\subsection*{You'll want to know: \hi{आवश्यक}}
\textbf{\hi{आवश्यक}} — \emph{āvaśyak} — \textquotedblleft{}essential, required (formal)\textquotedblright{}.

\begin{quote}
\hi{अगर} \hi{बीमारी} \hi{गंभीर} \hi{है} \hi{तो} \hi{आराम} \hi{आवश्यक} \hi{है।}
\end{quote}

\begin{quote}
\emph{agar bīmārī gambhīr hai to ārām āvaśyak hai.}
\end{quote}

\begin{quote}
If the illness is serious, rest is essential.
\end{quote}

\begin{cousinweb}
\textbf{\hi{आवश्यक}} (\emph{āvaśyak}) is a tatsama, borrowed from Sanskrit \emph{āvaśyaka}, \textquotedblleft{}necessary, inevitable\textquotedblright{} (Monier-Williams).
\end{cousinweb}

\begin{grammarlens}[title={What you've built}]
One more word for saying what happens if.
\end{grammarlens}

\subsection*{Your turn}
\begin{itemize}
\raggedright
  \item \emph{Say it:} \emph{āvaśyak}
  \item \emph{Say it:} \emph{āvaśyak}, once more
  \item \emph{Say it:} \emph{agar bīmārī gambhīr hai to ārām āvaśyak hai.} — If the illness is serious, rest is essential.
  \item \emph{Recall:} say the Hindi for a hardship, then the Hindi for to trap, then say \emph{āvaśyak} again
\end{itemize}

\begin{tcolorbox}[breakable,colback=teal!4,colframe=teal!35!black,title={Before you move on}]
What does \hi{आवश्यक} mean? (\textquotedblleft{}Essential, required (formal)\textquotedblright{}.) Say it once more.
\end{tcolorbox}
\section[\texorpdfstring{\hi{अनिवार्य} (anivārya)}{अनिवार्य (anivārya)}]{\hi{अनिवार्य} (anivārya) — compulsory}

\label{lesson:HI-C378-anivarya}



\begin{quote}
Before the new one: say the Hindi for to trap, then the Hindi for essential.
\end{quote}

\subsection*{You'll want to know: \hi{अनिवार्य}}
\textbf{\hi{अनिवार्य}} — \emph{anivārya} — \textquotedblleft{}compulsory\textquotedblright{}.

\begin{quote}
\hi{अगर} \hi{तुम} \hi{पहाड़} \hi{पर} \hi{चढ़ोगे} \hi{तो} \hi{रस्सी} \hi{अनिवार्य} \hi{है।}
\end{quote}

\begin{quote}
\emph{agar tum pahāṛ par caṛhoge to rassī anivārya hai.}
\end{quote}

\begin{quote}
If you climb the mountain, a rope is compulsory.
\end{quote}

\begin{cousinweb}
\textbf{\hi{अनिवार्य}} (\emph{anivārya}) is a tatsama, borrowed from Sanskrit \emph{anivārya}, \textquotedblleft{}not to be warded off, unavoidable\textquotedblright{} (Monier-Williams).
\end{cousinweb}

\begin{grammarlens}[title={What you've built}]
One more word for saying what happens if.
\end{grammarlens}

\subsection*{Your turn}
\begin{itemize}
\raggedright
  \item \emph{Say it:} \emph{anivārya}
  \item \emph{Say it:} \emph{anivārya}, once more
  \item \emph{Say it:} \emph{agar tum pahāṛ par caṛhoge to rassī anivārya hai.} — If you climb the mountain, a rope is compulsory.
  \item \emph{Recall:} say the Hindi for to trap, then the Hindi for essential, then say \emph{anivārya} again
\end{itemize}

\begin{tcolorbox}[breakable,colback=teal!4,colframe=teal!35!black,title={Before you move on}]
What does \hi{अनिवार्य} mean? (\textquotedblleft{}Compulsory\textquotedblright{}.) Say it once more.
\end{tcolorbox}
\section[\texorpdfstring{\hi{केवल} (keval)}{केवल (keval)}]{\hi{केवल} (keval) — only, solely}

\label{lesson:HI-C378-keval}



\begin{quote}
Before the new one: say the Hindi for essential, then the Hindi for compulsory.
\end{quote}

\subsection*{You'll want to know: \hi{केवल}}
\textbf{\hi{केवल}} — \emph{keval} — \textquotedblleft{}only, solely\textquotedblright{}.

\begin{quote}
\hi{अगर} \hi{तुम} \hi{केवल} \hi{सच} \hi{बोलोगे} \hi{तो} \hi{डरने} \hi{की} \hi{कोई} \hi{बात} \hi{नहीं} \hi{होगी।}
\end{quote}

\begin{quote}
\emph{agar tum keval sac bologe to ḍarne kī koī bāt nahīṃ hogī.}
\end{quote}

\begin{quote}
If you tell nothing but the truth, there will be nothing to fear.
\end{quote}

\begin{cousinweb}
\textbf{\hi{केवल}} (\emph{keval}) is a tatsama, borrowed from Sanskrit \emph{kevala}, \textquotedblleft{}alone, only, mere\textquotedblright{} (Monier-Williams).
\end{cousinweb}

\begin{grammarlens}[title={What you've built}]
One more word for saying what happens if.
\end{grammarlens}

\subsection*{Your turn}
\begin{itemize}
\raggedright
  \item \emph{Say it:} \emph{keval}
  \item \emph{Say it:} \emph{keval}, once more
  \item \emph{Say it:} \emph{agar tum keval sac bologe to ḍarne kī koī bāt nahīṃ hogī.} — If you tell nothing but the truth, there will be nothing to fear.
  \item \emph{Recall:} say the Hindi for essential, then the Hindi for compulsory, then say \emph{keval} again
\end{itemize}

\begin{tcolorbox}[breakable,colback=teal!4,colframe=teal!35!black,title={Before you move on}]
What does \hi{केवल} mean? (\textquotedblleft{}Only, solely\textquotedblright{}.) Say it once more.
\end{tcolorbox}
\section[\texorpdfstring{\hi{सिवाय} (sivāy)}{सिवाय (sivāy)}]{\hi{सिवाय} (sivāy) — except (\hi{के} \hi{सिवाय})}

\label{lesson:HI-C378-sivay}



\begin{quote}
Before the new one: say the Hindi for compulsory, then the Hindi for solely.
\end{quote}

\subsection*{You'll want to know: \hi{सिवाय}}
\textbf{\hi{सिवाय}} — \emph{sivāy} — \textquotedblleft{}except (\hi{के} \hi{सिवाय})\textquotedblright{}.

\begin{quote}
\hi{अगर} \hi{तुम्हारे} \hi{सिवाय} \hi{कोई} \hi{नहीं} \hi{आएगा} \hi{तो} \hi{हम} \hi{क्या} \hi{करेंगे}?
\end{quote}

\begin{quote}
\emph{agar tumhāre sivāy koī nahīṃ āegā to ham kyā kareṅge?}
\end{quote}

\begin{quote}
If nobody but you comes, what will we do?
\end{quote}

\begin{cousinweb}
\textbf{\hi{सिवाय}} (\emph{sivāy}) is usually traced to Arabic \emph{siwā}, \textquotedblleft{}other than\textquotedblright{}, which reached Hindi through Persian.
\end{cousinweb}

\begin{grammarlens}[title={What you've built}]
One more word for saying what happens if.
\end{grammarlens}

\subsection*{Your turn}
\begin{itemize}
\raggedright
  \item \emph{Say it:} \emph{sivāy}
  \item \emph{Say it:} \emph{sivāy}, once more
  \item \emph{Say it:} \emph{agar tumhāre sivāy koī nahīṃ āegā to ham kyā kareṅge?} — If nobody but you comes, what will we do?
  \item \emph{Recall:} say the Hindi for compulsory, then the Hindi for solely, then say \emph{sivāy} again
\end{itemize}

\begin{tcolorbox}[breakable,colback=teal!4,colframe=teal!35!black,title={Before you move on}]
What does \hi{सिवाय} mean? (\textquotedblleft{}Except (\hi{के} \hi{सिवाय})\textquotedblright{}.) Say it once more.
\end{tcolorbox}
\section[\texorpdfstring{\hi{मना} \hi{करना} (manā karnā)}{मना करना (manā karnā)}]{\hi{मना} \hi{करना} (manā karnā) — to forbid; to say no}

\label{lesson:HI-C378-manakarna}



\begin{quote}
Before the new one: say the Hindi for solely, then the Hindi for except.
\end{quote}

\subsection*{You'll want to know: \hi{मना} \hi{करना}}
\textbf{\hi{मना} \hi{करना}} — \emph{manā karnā} — \textquotedblleft{}to forbid; to say no\textquotedblright{}.

\begin{quote}
\hi{अगर} \hi{डाक्टर} \hi{मना} \hi{करेगा} \hi{तो} \hi{मैं} \hi{मिठाई} \hi{नहीं} \hi{खाऊँगा।}
\end{quote}

\begin{quote}
\emph{agar ḍākṭar manā karegā to maiṁ miṭhāī nahīṃ khāūṅgā.}
\end{quote}

\begin{quote}
If the doctor forbids it, I won't eat sweets.
\end{quote}

Past, with \hi{ने} (he or she …): \hi{उसने} \hi{मना} \hi{किया}, \emph{usne manā kiyā}.

\begin{cousinweb}
\textbf{\hi{मना} \hi{करना}} (\emph{manā karnā}) joins \hi{मना} to \hi{करना}, \textquotedblleft{}to do\textquotedblright{}. \hi{मना} is usually traced to Arabic \emph{manʿ}, \textquotedblleft{}prevention, holding back\textquotedblright{}, which reached Hindi through Persian.
\end{cousinweb}

\begin{grammarlens}[title={What you've built}]
One more word for saying what happens if.
\end{grammarlens}

\subsection*{Your turn}
\begin{itemize}
\raggedright
  \item \emph{Say it:} \emph{manā karnā}
  \item \emph{Say it:} \emph{manā karnā}, once more
  \item \emph{Say it:} \emph{agar ḍākṭar manā karegā to maiṁ miṭhāī nahīṃ khāūṅgā.} — If the doctor forbids it, I won't eat sweets.
  \item \emph{Recall:} say the Hindi for solely, then the Hindi for except, then say \emph{manā karnā} again
\end{itemize}

\begin{tcolorbox}[breakable,colback=teal!4,colframe=teal!35!black,title={Before you move on}]
What does \hi{मना} \hi{करना} mean? (\textquotedblleft{}To forbid; to say no\textquotedblright{}.) Say it once more.
\end{tcolorbox}
